% Extensión estimada: 5--5.5 páginas; cuatro subsecciones y tres tablas.
% Fuentes: HAM10000 (metadatos, Codigo/data/raw); y_true del test en Codigo/outputs/results/swin_tiny_sweep_t06_les_cwnone_predictions.json;
% Codigo/docs/ideas/criterio_seleccion_plan_20260917.md; Codigo/docs/resul/split_lesion_resul_20260926.md (§4: con el
% criterio anterior, ConvNeXt t06 tenía sens. macro 0.700 frente a 0.691 de t02);
% Codigo/docs/resul/revision_experimentacion_20261003.md (§2.3, §2.8 y §5); src/evaluation/clinical_ner.py;
% scripts/melanoma_detection.py; scripts/lesion_split_report.py; scripts/statistical_tests.py (bootstrap_ci,
% paired_bootstrap_diff, mcnemar_exact); scripts/ensemble_bootstrap.py (paired_bootstrap_diff_flags).
% Revisión del autor del 18/10/2026: sin fechas ni historia del proceso; cada métrica con su justificación; las
% pruebas estadísticas vienen de la antigua Sección 3.9; la especificidad de 0.85 se justifica como decisión de diseño.
% Revisión del tutor del 21/10/2026: pasa después de los componentes y del pipeline (antes era 3.3); recibe las
% referencias de comparación del pipeline (fuentes de la cascada, oracle, línea base sin generador) que estaban en 3.7.

\section{Protocolo y criterios de evaluación}
\label{sec:met_protocolo}

Con los componentes y el \emph{pipeline} definidos, falta establecer cómo se evalúan. La evaluación sigue tres reglas
generales. En primer lugar, toda decisión sobre un modelo, como su configuración o un umbral, se toma con el conjunto de validación, y el conjunto
de prueba, que no interviene en esas decisiones, se usa solo para informar el resultado final. En segundo lugar, cada resultado se lee frente a referencias de
comparación que dicen cuánto vale. Finalmente, una diferencia entre dos modelos se afirma solo si la sostiene una
comparación estadística. Dentro de ese marco, la clasificación y la generación de texto no se miden igual, así que cada
tarea justifica sus métricas por separado.

\subsection{Métricas de clasificación}
\label{sec:met_seleccion}

Las métricas de clasificación tienen que reflejar que los errores entre clases no cuestan lo mismo. Clasificar \gls{mel} como \gls{nv} retrasa la extirpación de una lesión cuya letalidad depende del estadio en que se diagnostica, mientras que clasificar \gls{nv} como \gls{mel} lleva a extirpar sin necesidad una lesión benigna, y los errores entre clases benignas tienen menos consecuencias~\cite{tschandl2018ham10000,nci2026melanoma}. Como además \gls{nv} reúne dos tercios
de \emph{HAM10000}, un buen resultado global puede esconder un mal desempeño en las clases que más importan. Por lo
tanto, las métricas tienen que evaluar cada clase por separado y darles a todas el mismo peso.

Con esa condición, el trabajo usa dos métricas que promedian las siete clases sin ponderarlas
(Sección~\ref{sec:marco_eval_clasif}). La sensibilidad macro (Ecuación~\ref{eq:marco_acc_bal}) es el resultado
principal, porque mide qué fracción de cada clase reconoce el modelo y es la métrica principal de los trabajos
revisados en la Sección~\ref{sec:estado_arte_diagnostico}, con los que así se puede comparar. El F1 macro
(Ecuación~\ref{eq:marco_f1}), en cambio, es el criterio con que se elige la configuración de cada modelo, porque
corrige lo que la sensibilidad macro no ve: un modelo que asigna muchas imágenes a las clases minoritarias sube su
sensibilidad aunque se equivoque en las demás, y el F1 macro lo penaliza a través de la precisión de cada clase. Las
dos se informan juntas, y cuando divergen, la diferencia muestra cuánto gana el modelo en detectar las clases a costa
de dar falsas alarmas.

La elección de un modelo se hace en validación, en dos niveles. Entre las configuraciones de un mismo modelo se elige
la de mayor F1 macro, porque en esa comparación tienen que pesar las clases minoritarias. Dentro de cada
entrenamiento, en cambio, la detención temprana y la época que se conserva se deciden por el F1 ponderado. La razón
es el tamaño de la validación: con 14 imágenes de DF y 18 de VASC (Tabla~\ref{tab:met_tamanos}), un solo acierto más
en DF sube su sensibilidad en 0.07, de modo que el F1 macro cambia de una época a otra por pocas imágenes y podría
detener el entrenamiento por azar. El F1 ponderado, dominado por las clases con más imágenes, varía menos. Su costo es
que la época conservada puede no ser la mejor para las clases minoritarias, y por eso las prioridades clínicas se
aplican en las demás decisiones: la configuración se elige por F1 macro, el resultado principal es la sensibilidad
macro y el melanoma tiene su propio umbral.

Una misma configuración tampoco da siempre el mismo resultado, porque las métricas de las clases minoritarias dependen
de pocas imágenes. Por eso, el resultado de cada clasificador supervisado, entrenado tres veces
(Sección~\ref{sec:met_supervisados_ensamble}), se informa como media $\pm$ desvío estándar muestral de esas tres
ejecuciones. Con tres ejecuciones, sin embargo, el desvío es una estimación gruesa, así que las comparaciones entre
modelos no se apoyan en él, sino en las comparaciones estadísticas de la Sección~\ref{sec:met_estadistica}, que se calculan
sobre las 1\,503 imágenes de prueba.

La detección de melanoma (T2) necesita un criterio propio, porque si se toma como melanoma la clase de mayor
probabilidad, el costo de perder un melanoma no interviene en la decisión. Por eso el umbral sobre la probabilidad de
melanoma se fija en validación como el de mayor sensibilidad entre los que mantienen una especificidad de al menos
0.850, y se aplica sin cambios al conjunto de prueba. Así se prioriza no perder melanomas, con un tope a las falsas alarmas: en
validación, a lo sumo el 15\,\% de los casos sin melanoma se señala como sospechoso. El \gls{auc}
(Ecuación~\ref{eq:marco_auc}), que no depende del umbral, se informa como resumen de la discriminación.

\subsection{Evaluación de la generación de reportes}
\label{sec:met_eval_generacion}

A diferencia de una clase, un reporte no tiene una única respuesta correcta, y por eso el texto generado se evalúa en
dos niveles, según haya o no un texto de referencia. En los dos se evalúa solo el reporte: los bloques de razonamiento
que el generador escribe entre etiquetas \texttt{\textless{}think\textgreater{}} se descartan antes de medir, incluido
el que queda abierto cuando el límite de tokens corta el texto. Se descartan porque el clínico lee el reporte y no el
razonamiento, y porque, si se midiera el texto completo, los diagnósticos que el razonamiento solo menciona como
hipótesis contarían como aciertos.

El Nivel~1 compara el reporte con un texto de referencia mediante BLEU-4, ROUGE-L, METEOR, CIDEr y BERTScore
(Sección~\ref{sec:marco_eval_texto}). Solo puede aplicarse a \emph{MM-Skin} y \emph{SkinCoT}, los conjuntos que tienen
ese texto, y se mide sobre sus particiones de prueba, de 1\,650 y 457 imágenes. De ellas se descartan los pares sin
referencia, es decir, las imágenes de \emph{MM-Skin} sin descripción y los textos de \emph{SkinCoT} que están enteros
dentro de bloques de razonamiento, de modo que quedan 1\,414 y 426 pares. En \emph{MM-Skin}, Mini-VLM tiene reportes
para una parte de esas imágenes, y los dos generadores se comparan sobre los 1\,022 pares que comparten. Además,
cuando el límite de tokens llega mientras el generador todavía razona, el reporte queda vacío. Esos reportes no entran
en las métricas, porque no hay texto que comparar, pero su proporción se informa junto a ellas.

Los valores del Nivel~1 se informan, pero no se usan para elegir ningún modelo ni ninguna configuración, porque los
textos de referencia no son diagnósticos clínicos de la lesión. Lo que mide la comparación depende de quién escribió
esos textos (Tabla~\ref{tab:met_conjuntos}). En \emph{MM-Skin} mide el parecido con la descripción de un libro de
texto, y en \emph{SkinCoT}, el parecido con un razonamiento que un modelo de lenguaje escribió conociendo el
diagnóstico.

El Nivel~2, en cambio, no necesita referencia y se aplica al conjunto de prueba de \emph{HAM10000}, con las dos métricas sin
referencia de la Sección~\ref{sec:marco_eval_texto}. La cobertura de entidades clínicas (\emph{clinical entity
recall}) se informa como complemento, porque mide cuánto contenido clínico describe el texto y no si es pertinente. La
consistencia diagnóstica (\emph{diagnostic consistency}, Ecuación~\ref{eq:marco_fd}) es la métrica principal de
generación, porque responde directamente la segunda pregunta de investigación: si el reporte afirma el diagnóstico
correcto de la imagen que recibe.

Para calcular la consistencia, el diagnóstico se reconoce con una búsqueda estricta, que solo acepta los términos
distintivos de cada clase, como «melanoma» o «basal cell», y exige palabra completa en las abreviaturas, como «BCC».
Una búsqueda que aceptara abreviaturas o términos genéricos, en cambio, acreditaría diagnósticos que el texto no
nombra, como «mel» dentro de «melanocytic» o la queratosis benigna por la palabra «keratosis» de «actinic keratosis».

La consistencia dice si el reporte nombra la clase verdadera, pero no cómo falla cuando no lo hace. Por eso cada
reporte se asigna además a su falla principal, según las categorías de la Tabla~\ref{tab:marco_fallas}, con reglas
sobre el texto y la misma búsqueda estricta. Nombrar dos o más clases es una contradicción, nombrar solo una clase
distinta de la verdadera es un diagnóstico incorrecto, y no nombrar ninguna, con expresiones de reserva como «may be»
o «cannot be determined», es \emph{hedging}. En cambio, una expresión de reserva junto a la clase verdadera, como
«consistent with nevus», no se cuenta como \emph{hedging}, porque el reporte se compromete con el diagnóstico.

Aun así, la consistencia diagnóstica mide fidelidad diagnóstica y no validez clínica, porque, como las demás métricas
de generación, no comprueba que lo descrito esté en la imagen (Sección~\ref{sec:marco_eval_texto}). Por eso ningún
resultado se presenta como evidencia de validez clínica. La Tabla~\ref{tab:met_roles} resume el papel de cada métrica
de clasificación y de generación en el trabajo.

\begin{table}[!tbp]
    \caption{Papel de cada métrica en el trabajo.}
    \label{tab:met_roles}
    \begin{tabularx}{\textwidth}{@{}>{\raggedright\arraybackslash}p{4.6cm}
        >{\raggedright\arraybackslash\hspace{0pt}}X >{\raggedright\arraybackslash\hspace{0pt}}p{3.6cm}@{}}
        \toprule
        \textbf{Métrica} & \textbf{Papel en el trabajo} & \textbf{Se define en} \\
        \midrule
        Sensibilidad macro & Resultado principal de clasificación en el conjunto de prueba & Sección~\ref{sec:marco_eval_clasif} \\
        \filasep
        F1 macro & Elige la configuración de cada modelo (validación); se informa en el conjunto de prueba & Ecuación~\ref{eq:marco_f1} \\
        \filasep
        F1 ponderado & Elige la época dentro de cada entrenamiento (validación) & Ecuación~\ref{eq:marco_f1} \\
        \filasep
        Exactitud & Se informa & Sección~\ref{sec:marco_eval_clasif} \\
        \filasep
        Sensibilidad y especificidad de melanoma & Umbral fijado en validación (especificidad de al menos 0.850); resultado en el conjunto de prueba & Sección~\ref{sec:marco_eval_clasif} \\
        \filasep
        AUC de melanoma & Se informa & Sección~\ref{sec:marco_eval_clasif} \\
        \filasep
        Nivel 1 (BLEU-4, ROUGE-L, METEOR, CIDEr, BERTScore) & Se informa, solo en \emph{MM-Skin} y \emph{SkinCoT} & Sección~\ref{sec:marco_eval_texto} \\
        \filasep
        Cobertura de entidades clínicas & Complemento del Nivel 2 & Sección~\ref{sec:marco_eval_texto} \\
        \filasep
        Consistencia diagnóstica & Métrica principal de generación (solo el reporte, búsqueda estricta) & Sección~\ref{sec:marco_eval_texto} \\
        \filasep
        Calibración (ECE, Brier) & Se analiza en la Sección~\ref{sec:met_calibracion} & Sección~\ref{sec:marco_eval_calib} \\
        \bottomrule
    \end{tabularx}
\end{table}

\subsection{Referencias de comparación}
\label{sec:met_referencias}

Un valor de una métrica no dice por sí solo si un modelo es bueno, y por eso cada resultado se lee frente a una
referencia. En la clasificación se usan tres. En primer lugar, el azar fija el piso de cualquier clasificador de
siete clases, una sensibilidad macro de 1/7. En segundo lugar, el clasificador trivial que responde siempre NV muestra
lo que se obtiene del desbalance sin mirar la imagen: en el conjunto de prueba alcanza una exactitud de 0.669 y un F1 ponderado de
0.537, pero una sensibilidad macro de 0.143 y un F1 macro de 0.115. Por eso la exactitud y el F1 ponderado, que
favorecen a NV, se informan pero no son el resultado principal. Finalmente, los clasificadores supervisados de la
Sección~\ref{sec:met_supervisados} son la línea base de los \glspl{vlm}, y la primera pregunta de investigación se
responde comparándolos.

En la generación, las referencias sirven para separar lo que aporta cada componente del \emph{pipeline}. Para ello, la
cascada de la Sección~\ref{sec:met_cascada} se ejecuta con el mismo generador y cuatro fuentes de diagnóstico:
\begin{itemize}[leftmargin=1.6em]
    \item \textbf{Swin-Tiny}, el clasificador del \emph{pipeline}.
    \item \textbf{DermLIP en modo zero-shot}, una alternativa sin entrenamiento sobre \emph{HAM10000}
        (Sección~\ref{sec:met_clasif_vlm}).
    \item \textbf{\emph{Oracle}}, que entrega el diagnóstico verdadero. Es una cota superior, es decir, el mejor
        resultado que alcanzaría la cascada si el clasificador no se equivocara, y no un sistema usable, porque en la
        práctica la etiqueta no se conoce de antemano.
    \item \textbf{La ausencia de anclaje}, es decir, el generador sin diagnóstico, que muestra qué aporta el anclaje al
        reporte.
\end{itemize}
En los cuatro casos se usan los tres generadores de la Sección~\ref{sec:met_generacion}: Mini-VLM y LLaVA-Med con los
pesos ajustados allí, y SkinVL con sus pesos publicados. Las dos plantillas de recuperación
(Sección~\ref{sec:met_recuperacion}) se comparan con la cascada sobre las mismas imágenes.

Además, dos variantes de la cascada con el \emph{oracle} separan lo que el generador toma de la imagen de lo que toma
del diagnóstico. En la primera, la imagen se reemplaza por una gris neutra, como en el control de la
Sección~\ref{sec:met_clasif_vlm}. En la segunda, el diagnóstico es equivocado: cada clase se cambia por otra con la
que se suele confundir, como \gls{mel} por \gls{nv}. Mini-VLM y LLaVA-Med se evalúan en las dos sobre un subconjunto
fijo de 400 imágenes del conjunto de prueba, por el tiempo que lleva cada reporte, y SkinVL, que escribe más rápido, sobre el conjunto de prueba completo.

Por último, el ruteo tiene su propia línea base, que aísla el aporte del generador. La línea base sin generador usa el
mismo ruteo y el mismo clasificador, pero no genera texto: el nivel afirmativo devuelve el diagnóstico del
clasificador, el diferencial devuelve el par de candidatos y el nivel de derivación queda igual que en el
\emph{pipeline}. Junto con el \emph{oracle}, que fija la cota superior, es la referencia con la que se compara el
\emph{pipeline} completo.

\subsection{Comparación estadística}
\label{sec:met_estadistica}

Una vez calculadas las métricas, que un modelo obtenga en el conjunto de prueba un valor de una métrica más alto que otro, o
viceversa, no alcanza para afirmar que uno sea mejor que el otro (Sección~\ref{sec:marco_eval_estadistica}). Por
ello, cada comparación usa el \emph{bootstrap} o el test de McNemar exacta, según lo que se pregunta. Todos los
\emph{bootstrap} usan 2\,000 remuestreos de las imágenes de prueba. La unidad de las comparaciones es la imagen de prueba
y no la ejecución, de modo que se calculan sobre una de las tres ejecuciones de cada clasificador, la
\textbf{ejecución de referencia}, que es también la que usan el \emph{pipeline} y las evaluaciones complementarias. La
variación entre ejecuciones se informa aparte como desvío.

Como se comparan siete clasificadores (los cuatro supervisados, DermLIP sin ajuste y los dos \glspl{vlm} generativos
sin ajuste, Sección~\ref{sec:met_clasif_vlm}), que forman 21 pares, los valores p de McNemar se corrigen con el método
de Holm (Ecuación~\ref{eq:marco_holm}). La sensibilidad de cada clase, en cambio, se compara de forma descriptiva,
contando en cuántas de las siete clases mejora cada modelo, porque con siete valores ningún test de rangos podría
alcanzar el umbral que Holm exige para 21 comparaciones. Además, el \emph{bootstrap} remuestrea imágenes y no
lesiones, así que trata como independientes las imágenes de una misma lesión. En el conjunto de prueba hay 1\,503 imágenes de
1\,429 lesiones, y remuestrear por lesión casi no cambia los intervalos. La Tabla~\ref{tab:met_pruebas} indica qué test se aplica a cada comparación.

\begingroup
\small
\setlength{\tabcolsep}{3pt}
\begin{table}[!tbp]
    \centering
    \caption{Comparaciones estadísticas según la comparación.}
    \label{tab:met_pruebas}
    \begin{tabularx}{\textwidth}{@{}>{\raggedright\arraybackslash}p{3.4cm}
        >{\raggedright\arraybackslash\hspace{0pt}}X
        >{\raggedright\arraybackslash\hspace{0pt}}X
        >{\raggedright\arraybackslash\hspace{0pt}}p{2.6cm}@{}}
        \toprule
        \textbf{Comparación} & \textbf{Pregunta que responde} & \textbf{Test} & \textbf{Criterio} \\
        \midrule
        Métrica agregada de un modelo
            & ¿Cuánto varía por la composición del conjunto de prueba?
            & Intervalo percentil de \emph{bootstrap}
            & IC del 95\,\% \\
        \filasep
        Diferencia de F1 macro o de sensibilidad macro entre dos modelos
            & ¿La diferencia excede la variación de la muestra?
            & \emph{Bootstrap} pareado
            & IC del 95\,\% sin el cero \\
        \filasep
        Aciertos por imagen entre dos modelos
            & ¿Uno acierta imágenes en las que el otro falla?
            & McNemar exacto sobre los casos discordantes, con corrección de Holm
            & $p < 0.05$ corregido \\
        \filasep
        Sensibilidad por clase entre dos modelos
            & ¿En cuántas clases mejora cada uno?
            & Descriptiva
            & --- \\
        \filasep
        Consistencia diagnóstica entre dos variantes del \emph{pipeline}
            & ¿La diferencia excede la variación de los reportes?
            & \emph{Bootstrap} pareado de la diferencia de proporciones
            & IC del 95\,\% sin el cero \\
        \bottomrule
    \end{tabularx}
\end{table}
\endgroup

Estos criterios se aplican a los resultados sobre el conjunto de prueba de \emph{HAM10000}, que reproduce las condiciones en que
se entrenaron los modelos. Las evaluaciones complementarias de la sección siguiente cambian esas condiciones.
